\section{Introduction}
\label{sec:introduction}

Deep Reinforcement Learning (DRL) has emerged as one of the most powerful
paradigms for learning complex control policies directly from raw sensory
data~\citep{mnih2015humanlevel,silver2016mastering}. By combining the
representational capacity of deep neural networks with the sequential
decision-making framework of Reinforcement Learning (RL), DRL agents have
achieved impressive results in domains ranging from board games and video
games to robotic manipulation and continuous locomotion~\citep{schulman2017proximal,haarnoja2018soft}.
A particularly appealing property of DRL is its ability to operate end-to-end:
given only raw observations and a scalar reward signal, an agent can learn
both a useful representation of the world and a policy for acting in it,
without access to hand-crafted features or an explicit model of the environment.

Among the earliest and most influential families of DRL algorithms are those
based on the \emph{action-value function}, or $Q$-function, which estimates
the expected discounted cumulative reward of taking a given action in a given
state and following the optimal policy thereafter. The seminal
Deep Q-Network (DQN)~\citep{mnih2015humanlevel} demonstrated that a
convolutional neural network (CNN) could approximate the optimal
$Q^{*}(s,a)$ directly from pixel observations, using two key stabilising
mechanisms: a \emph{replay buffer} that decorrelates temporally adjacent
transitions before they are used for gradient updates, and a periodically
frozen \emph{target network} that provides stable regression targets and
prevents the harmful feedback loop between the online network and its own
bootstrapped targets. Despite its relative simplicity, DQN established
a compelling proof of concept and remains an important baseline for
value-based methods.

A prominent milestone in this line of research is \textbf{Rainbow
DQN}~\citep{hessel2018rainbow}, which integrates six orthogonal improvements
over the original algorithm into a single, unified agent. These are:
Double DQN~\citep{van2016deep}, which decouples action selection from action
evaluation to reduce overestimation bias;
Prioritised Experience Replay~\citep{schaul2016prioritized}, which samples
transitions proportionally to their temporal-difference (TD) error so that
more informative experiences are replayed more often;
the Dueling Network architecture~\citep{wang2016dueling}, which decomposes
$Q(s,a)$ into a state-value stream $V(s)$ and an advantage stream $A(s,a)$
to improve generalisation across actions;
multi-step returns~\citep{sutton2018reinforcement}, which replace one-step
TD targets with $n$-step Monte Carlo estimates that propagate reward
information more rapidly;
Distributional RL~\citep{bellemare2017distributional}, which models the full
return distribution rather than its expectation; and
Noisy Networks~\citep{fortunato2018noisy}, which replace $\varepsilon$-greedy
exploration with learnable parameter noise.
Together, these components yield substantially higher sample efficiency and
final performance compared with either DQN or any single extension applied
in isolation, making Rainbow a natural next step after the DQN baseline.

\subsection{Motivation and Scope}

A fundamental challenge in applying value-based DRL to physically realistic
continuous control is the \emph{mismatch between continuous action spaces and
discrete $Q$-function outputs}. Environments provided by
MuJoCo~\citep{todorov2012mujoco} via the DeepMind Control
Suite (dm\_control)~\citep{tunyasuvunakool2020dmcontrol} expose
high-dimensional, continuous action vectors—torques applied simultaneously to
multiple articulated joints—which are not directly compatible with
$Q$-function methods that enumerate actions at the output layer.
A carefully designed \emph{action discretisation} strategy is therefore
required to bridge this gap while preserving enough behavioural richness for
the agent to learn coordinated locomotion gaits.

A second challenge is learning from \emph{pixel observations} alone.
Rather than providing the agent with privileged state vectors
(joint angles, velocities, contact forces), in this work the agent receives
only raw RGB images of the simulated scene. The network must therefore
simultaneously learn a compact and informative visual representation of the
environment and a high-quality value function over the discretised action
space—a dual objective that substantially increases the difficulty and
computational cost of training.

A third challenge is \emph{task complexity and morphological difficulty}.
Legged locomotion spans a wide spectrum: from low-dimensional systems with
inherently stable gaits to high-dimensional humanoid bodies whose dynamics
are notoriously difficult to control. In this work we study both ends of this
spectrum. For the DQN baseline we use the \textbf{Hopper} environment, a
one-legged planar system with three controllable joints, whose relatively
low dimensionality makes it a tractable testbed for understanding the
limitations of standard DQN under pixel observations. For the Rainbow
agent we use the \textbf{Humanoid} environment, a $21$-degrees-of-freedom
full-body model whose high control dimensionality, unstable balance, and
complex contact dynamics make it one of the most challenging continuous
control benchmarks available, allowing us to assess whether Rainbow's
combined improvements are sufficient to tackle substantially harder tasks.

\subsection{Contributions and Structure}

In this report we present a comparative experimental study of value-based DRL
for pixel-based locomotion in MuJoCo environments of increasing difficulty.
Specifically, our contributions are as follows:

% TODO: Change DQN to include branching DQN?
\begin{enumerate}
  \item We implement and train a standard DQN agent with a CNN encoder on the
        Hopper locomotion task under pixel observations, analysing the effect
        of action discretisation and reward design on learning stability and
        final performance.
  \item We implement Rainbow DQN and apply it to the significantly harder
        Humanoid locomotion task, studying which of Rainbow's six components
        are most critical for learning in a high-dimensional, unstable control
        setting.
  \item We provide a direct qualitative and quantitative comparison between
        both agents across training dynamics and evaluation performance,
        discussing the extent to which Rainbow's improvements translate to
        gains in sample efficiency and policy quality under the additional
        difficulty of pixel-only observations.
\end{enumerate}

% TODO: put correct section references
The remainder of this report is organised as follows.
Section~\ref{sec:dqn} details the DQN baseline on Hopper: environment
configuration, action discretisation, reward shaping, network architecture,
and training results.
Section~\ref{sec:rainbow} presents the Rainbow DQN agent on Humanoid,
including implementation details and a comparison against the DQN baseline.
Section~\ref{sec:conclusions} draws general conclusions, reflects on the
lessons learned across both experimental settings, and outlines directions
for future work.